\documentclass[12pt, a4paper, twocolumn]{article}

% ===================================================================
% PAQUETES ESENCIALES
% ===================================================================

% Codificación y lenguaje
\usepackage[utf8]{inputenc}
\usepackage[spanish]{babel}

% Matemáticas
\usepackage{amsmath}
\usepackage{amssymb}
\usepackage{amsfonts}

% Gráficos e imágenes
\usepackage{graphicx}
\graphicspath{{../assets/images/}}

% Configuración de página
\usepackage{geometry}
\geometry{a4paper, margin=0.75in}

% Enlaces e hipervínculos
\usepackage{hyperref}
\hypersetup{
    colorlinks=true,
    linkcolor=blue,
    filecolor=magenta,
    urlcolor=cyan,
    pdftitle={Título del Artículo},
    pdfauthor={Autor},
}

% Resumen en español
\usepackage{abstract}

% ===================================================================
% INFORMACIÓN DEL DOCUMENTO
% ===================================================================

\title{\textbf{Título del Artículo Científico}}
\author{
    Nombre del Autor\textsuperscript{1} \\
    \small \textsuperscript{1}Institución/Universidad \\
    \small \texttt{email@ejemplo.com}
}
\date{\today}

% ===================================================================
% INICIO DEL DOCUMENTO
% ===================================================================

\begin{document}

\maketitle

\begin{abstract}
Escribe aquí el resumen del artículo. Debe ser conciso y describir el objetivo, metodología, resultados principales y conclusiones del trabajo.
\end{abstract}

\textbf{Palabras clave:} palabra1, palabra2, palabra3

\section{Introducción}

La introducción debe presentar el contexto del trabajo, la problemática a abordar y los objetivos del estudio.

\section{Metodología}

Describe los métodos utilizados en el estudio.

\section{Resultados}

Presenta los resultados obtenidos.

\subsection{Análisis de datos}

Análisis detallado de los datos.

\section{Discusión}

Discute los resultados en el contexto de la literatura existente.

\section{Conclusiones}

Resume las conclusiones principales del trabajo.

\section*{Agradecimientos}

Agradecimientos opcionales.

\begin{thebibliography}{9}
\bibitem{referencia1}
Apellido, N. (Año). \textit{Título del artículo}. Revista, Volumen(Número), páginas.

\bibitem{referencia2}
Apellido, N. (Año). \textit{Título del libro}. Editorial.
\end{thebibliography}

\end{document}
